\preface

Scattering amplitudes have been termed the ``most perfect microscopic structures in the Universe''.\footnote{L.~J.~Dixon,
{\it ``Scattering amplitudes: the most perfect microscopic structures in the universe,''}
J. Phys. A \textbf{44} (2011), 454001.} 
Scattering amplitudes are elementary building blocks in quantum field theory that allow us to predict probabilities for the outcome of particle collisions. 
Therefore they are a crucial link connecting theory and experiment. For example, they allow us to test predictions from the Standard Model of particle physics against the collider data being collected at the Large Hadron Collider (LHC) at CERN. 


\smallskip
The study of scattering amplitudes in quantum field theory has a long history,
dating back to the analytic $S$-matrix program of the 1960's. The modern field can trace its
origins to the 1980's. At that time, the state of the art of amplitude
computation was five-gluon scattering at tree level, i.e.\ the lowest order in
perturbation theory. Parke and Taylor famously managed to simplify what were
page-long results to a beautiful single-line formula, thus providing the first
of many hints for the underlying simplicity of scattering amplitudes in gauge
theory. To date, the state of the art has advanced by several loop orders,
i.e.\ to higher orders in the perturbative expansion in the coupling constant.
This was made possible by many conceptual advances. Modern approaches often use
gauge-invariant building blocks, as opposed to the traditional Feynman
diagrams, to organise calculations.
%
Perturbative unitarity made it possible to write down relations that recycle
these gauge-invariant building blocks, effectively leading to a recursion, both
in the number of particles and in the loop order. This is closely related to a
deeper understanding of the analytic structure of scattering amplitudes at loop level.

 \smallskip 
 
Our lecture notes provide an introduction for anyone wishing to learn more about this fascinating subject.
They are suitable for students at M.Sc.\ or Ph.D.\ levels.
We have made a specific selection of topics, so that they can be used for a one-semester university lecture series. This is based on our teaching experience, both at M.Sc.\ level at the universities of Berlin and Turin, as well as at Ph.D.-level summer schools worldwide. Additionally, they may be a useful primer for graduate students wishing to do research in this area. 
The chapters of the book cover necessary quantum field theory background, on-shell techniques for tree-level and one-loop scattering amplitudes, as well as dedicated techniques for evaluating Feynman loop integrals. 
Many exercises complement the text, and we provide fully worked-out solutions, as well as examples of how to implement the material using computer algebra codes. 
Supplementary material, \textsc{Mathematica} notebooks, corrections and further information are provided and maintained at the following 
dedicated website:
\begin{center}
\href{https://scattering-amplitudes.mpp.mpg.de/scattering-amplitudes-in-qft/}{https://scattering-amplitudes.mpp.mpg.de/scattering-amplitudes-in-qft/} \,.
\end{center}


\smallskip
Our aim is to provide a useful starting point to enable readers to access and contribute to current research in this thriving field.
Indeed, all the advancements mentioned above have significantly contributed to the community's ability to make more accurate predictions for particle scattering processes relevant to collider physics.
Beyond that, recent applications use these multi-loop techniques in Einstein's theory of gravity to provide efficient, high-precision predictions for the gravitational waveforms emitted in the encounter of black holes and
neutron stars. Therefore they provide the basis for data analysis in present and future gravitational wave detectors.
Furthermore, some of the properties that were revealed in the recent studies of scattering amplitudes in quantum field theories, 
such as hidden symmetries, dualities, and conceptually new approaches, 
have taught us entirely novel ways of thinking about quantum field theory. 
This may ultimately lead to a reformulation of the theory, where Feynman diagrams no longer play a fundamental role.
%\smallskip
Since 2009, progress in this fast-growing field of theoretical physics is evident via the ``Amplitudes'' conferences held annually, which bring together researchers 
interested in both formal and phenomenological aspects of scattering amplitudes, with a wide range of applications from pure mathematics to collider and gravitational wave physics. 


\smallskip

Some elements of modern on-shell methods for scattering amplitudes are also discussed in the quantum field theory textbooks
of Srednicki~{\it ``Quantum field theory''},
  Zee~{\it ``Quantum field theory in a nutshell''}
  and Schwartz~{\it ``Quantum Field Theory and the Standard Model''}, as well as in the dedicated reviews by Mangano and Parke~{\it ``Multiparton amplitudes in gauge theories''}, and by Dixon~{\it ``Calculating scattering amplitudes efficiently''}.
Other recent reviews on topics covered in these lecture notes include those by Weinzierl~{\it ``Tales of 1001 Gluons''}, and {\it ``Feynman Integrals''}. 
These offer a comprehensive account which can be especially interesting to mathematically-inclined readers.
For advanced topics in scattering amplitudes, focusing in particular on developments in supersymmetric theories, we highly recommend the textbook by Elvang and Huang~{\it ``Scattering Amplitudes in Gauge Theory and Gravity''}
 as well as the research monograph 
by Arkani-Hamed, Bourjaily, Cachazo, Goncharov, Postnikov and Trnka~{\it ``Grassmannian Geometry of Scattering Amplitudes''}.



\vspace{\baselineskip}
\begin{flushright}\noindent
Berlin, Munich and Turin\hfill {\it Simon Badger}\\
June 2023\hfill {\it Johannes M.\ Henn}\\
{\it Jan C.\ Plefka}\\
{\it Simone Zoia}\\
\end{flushright}